\documentclass[17pt]{extarticle}  % print quality
\usepackage{polyglossia}
\usepackage[margin=20mm]{geometry}
\setdefaultlanguage{marathi}
\setotherlanguages{english}
\usepackage{fourier}
%\textbf with "Sanskrit Text" produces the shape warning because the default Sanskrit Text font on Windows does not have a bold equivalent; try Yashomudra or Shobhika; now I like Tiro better
\newfontfamily\devanagarifont[Script=Devanagari]{Sanskrit Text}
%\newfontfamily\devanagarifont[Script=Devanagari]{Tiro Devanagari Marathi}

\pagenumbering{gobble}
\setlength{\parindent}{0pt}% Remove paragraph indent
\usepackage[skip=\medskipamount]{parskip}
\usepackage{fontawesome}
\usepackage[style=english]{csquotes}
\usepackage{verse}
\usepackage{xcolor}
\usepackage{hyperref}
\hypersetup{
    colorlinks=true,
    linkcolor=blue,
    filecolor=magenta,
    citecolor=blue,
    urlcolor=purple,
}

% control hyphenation: https://tex.stackexchange.com/a/177179/64425
\tolerance=1
\emergencystretch=\maxdimen
\hyphenpenalty=10000
\hbadness=10000
% control hyphenation: https://tex.stackexchange.com/a/177179/64425

\begin{document}

\title{एक धागा सुखाचा (जगाच्या पाठीवर)}
\author{ग. दि. माडगूळकर}
\date{}
\maketitle
\hrule
\vspace{5mm}

%\textenglish {

%This document is a sample document to test font families and font typefaces.
%}

एक धागा सुखाचा, शंभर धागे दु:खाचे\\
जरतारी हें वस्त्र माणसा तुझिया आयुष्याचें
\vspace{\baselineskip}

\vin पांघरसी जरि असला कपडा\\
\vin येसी उघडा, जासी उघडा\\
कपड्यासाठीं करिसी नाटक, तीन प्रवेशांचें !
\vspace{\baselineskip}

\vin मुकीं अंगडीं बालपणाचीं\\
\vin रंगित वसनें तारुण्याचीं\\
जीर्ण शाल मग उरे शेवटीं, लेणें वार्धक्याचें !
\vspace{\baselineskip}

\vin या वस्त्रांतें विणतो कोण ?\\
\vin एकसारखीं नसतीं दोन\\
कुणा न दिसले त्रिखंडांत, त्या हात विणकर्‍याचे !



\vspace{5mm}
\hrule
\end{document}
